\documentclass{article}
\usepackage{tekvideoexercises}

\begin{document}
\exercisename{Definition af logaritmer}
\streaklength{5}

\tableofcontents
\newpage

\begin{exercise}{Definition af logaritmer 1-1}

Bestem $x$ således at $2^x = 8$ og skriv svaret som en logaritme.

\answerbox{$x = \log_2(8)$}

\hint

Definitionen af en logaritme er: Hvis $a^x = b$, så er $x = \log_a(b)$.

\hint

Anvend definitionen direkte:

\hint
\[
2^x = 8 \implies x = \log_2(8)
\]

\end{exercise}

\newpage

\begin{exercise}{Definition af logaritmer 1-2}

Bestem $y$ således at $3^y = 27$ og skriv svaret som en logaritme.

\answerbox{$y = \log_3(27)$}

\hint

Brug definitionen af logaritmer:

\hint
\[
3^y = 27 \implies y = \log_3(27)
\]

\end{exercise}

\newpage

\begin{exercise}{Definition af logaritmer 1-3}

Bestem $x$ således at $5^x = 1$ og skriv svaret som en logaritme.

\answerbox{$x = \log_5(1)$}

\hint

Brug definitionen:

\hint
\[
5^x = 1 \implies x = \log_5(1)
\]

\hint

Bemærk at $\log_5(1) = 0$ fordi $5^0 = 1$.

\end{exercise}

\newpage

\begin{exercise}{Definition af logaritmer 1-4}

Bestem $x$ således at $10^x = 100$ og skriv svaret som en logaritme.

\answerbox{$x = \log_{10}(100)$}

\hint

Brug definitionen:

\hint
\[
10^x = 100 \implies x = \log_{10}(100)
\]

\hint

Bemærk at $\log_{10}(100) = 2$ fordi $10^2 = 100$.

\end{exercise}

\newpage

\begin{exercise}{Definition af logaritmer 1-5}

Bestem $x$ således at $b^x = a$ (hvor $a > 0$ og $b > 0$, $b \neq 1$) og skriv svaret som en logaritme.

\answerbox{$x = \log_b(a)$}

\hint

Dette er den generelle definition af en logaritme.

\hint
\[
b^x = a \implies x = \log_b(a)
\]

\end{exercise}

\newpage

\begin{exercise}{Definition af logaritmer 1-6}

Bestem $x$ således at $4^x = 64$ og skriv svaret som en logaritme.

\answerbox{$x = \log_4(64)$}

\hint

Brug definitionen:

\hint
\[
4^x = 64 \implies x = \log_4(64)
\]

\hint

Bemærk at $\log_4(64) = 3$ fordi $4^3 = 64$.

\end{exercise}

\end{document}
